\section{统一求解与适用边界}

\subsection{状态判别与数值闭合}

三问共用浮标、刚性构件和锚链静力模型。给定环境与设计变量后，程序由下向上递推竖直张力和构件倾角，再用水深条件闭合吃水、锚链高度与各段竖直投影。锚链先按部分铺底分支求解；若所需悬起长度超过总链长，则切换到完全悬起分支并引入锚点竖直张力。该处理避免预先固定锚链状态造成不相容解。

问题二在静力求解器外对配重质量作区间搜索，分别求钢桶倾角和锚点夹角的临界质量并取较大者；问题三枚举链型、链环数和配重档位，再对候选方案作多工况复核。每个结果均检查水深闭合、受力平衡、吃水范围和锚链状态，问题一另以210链环离散模型校核连续悬链线近似。

\subsection{模型边界}

统一模型采用准静态、不可伸长柔索和准定常阻力近似，不能反映结构惯性、链环弹性、海床接触滞回及风流脉动造成的瞬态响应。状态切换由解析边界判定，在临界离床区域可能放大数值不光滑性；代表工况枚举也只验证离散环境集合。因而本文结论用于稳态方案比较，不能替代时域动力学、接触力学和连续不确定环境下的鲁棒性分析。
